\chapter{Равномерная сходимость функциональных последовательностей и рядов.
Непрерывность, интегрируемость и дифференцируемость суммы функционального
ряда}

\section{Функциональные последовательности}

\begin{thm}[Непрерывность предельной функции] \label{ch16:thm:cont}
Пусть функции $f_n$, $n = 1, 2, \ldots$, непрерывны на множестве $E$, причём $f_n(x) \overset{E}{\rightrightarrows} f(x)$. Тогда предельная функция $f$ непрерывна на~$E$.
\end{thm}

\begin{proof}
Из равномерной сходимости следует, что
\begin{equation} \label{ch16:eq:unif}
\forall\, \varepsilon > 0 \;\; \exists\, n_0 \cquad \forall\, n \ge n_0, \;\; \forall\, x \in E \cquad |f_n(x) - f(x)| < \frac{\varepsilon}{3}.
\end{equation}

Пусть $x_0, x \in E$. Оценим приращение предельной функции:
\begin{equation} \label{ch16:eq:triangle}
|f(x) - f(x_0)| \le |f(x) - f_n(x)| + |f_n(x) - f_n(x_0)| + |f_n(x_0) - f(x_0)|.
\end{equation}

В силу~\eqref{ch16:eq:unif} первое и третье слагаемые в правой части~\eqref{ch16:eq:triangle} меньше $\dfrac{\varepsilon}{3}$ при $n \ge n_0$.

Зафиксируем $x_0 \in E$ и $n \ge n_0$. Так как функция $f_n$ непрерывна на множестве $E$, то она непрерывна в точке $x_0$ по множеству $E$:
$$
\forall\, \varepsilon > 0 \;\; \exists\, \delta > 0 \cquad \forall\, x \in U_\delta(x_0) \cap E \cquad |f_n(x) - f_n(x_0)| < \frac{\varepsilon}{3}.
$$

Из~\eqref{ch16:eq:triangle} окончательно получаем:
$$
\forall\, \varepsilon > 0 \;\; \exists\, \delta > 0 \cquad \forall\, x \in U_\delta(x_0) \cap E \cquad |f(x) - f(x_0)| < \frac{\varepsilon}{3} + \frac{\varepsilon}{3} + \frac{\varepsilon}{3} = \varepsilon.
$$

Значит, $f$ непрерывна в точке $x_0$ по множеству $E$. Так как $x_0 \in E$ "--- произвольная точка, $f$ непрерывна на~$E$.
\end{proof}

\begin{notion}
В случае поточечной (но неравномерной) сходимости последовательности непрерывных функций предельная функция не обязана быть непрерывной. Разрывность предельной функции является косвенным доказательством неравномерности сходимости.

Вместе с тем равномерная сходимость последовательности непрерывных функций является \textit{достаточным}, но не необходимым условием непрерывности предельной функции.
\end{notion}

\begin{notion}
Теорема~\ref{ch16:thm:cont} справедлива при $E \subset \bbR$, $E \subset \bbR^n$, $E \subset \bbC$. В последнем случае можно рассматривать комплекснозначные функции комплексной переменной; определения предела и непрерывности ничем не отличаются от действительного случая (модуль "--- модуль комплексного числа, $U_\delta(x_0) = \{x \in \bbC \colon |x - x_0| < \delta\}$).

Теоремы об интегрировании и дифференцировании доказываются в действительном случае для отрезка.
\end{notion}

\section{Теорема об интегрировании}

\begin{thm}[Почленное интегрирование] \label{ch16:thm:int}
Пусть функции $f_n$ непрерывны на отрезке $[a;\, b]$, причём $f_n(x) \rightrightarrows f(x)$, $x \in [a;\, b]$. Тогда
$$
\int_a^x f_n(t)\,dt \rightrightarrows \int_a^x f(t)\,dt, \quad x \in [a;\, b].
$$
В частности, при $x = b$:
$$
\int_a^b f_n(t)\,dt \to \int_a^b f(t)\,dt
$$
"--- предельный переход под знаком интеграла.
\end{thm}

\begin{proof}
По теореме~\ref{ch16:thm:cont} функция $f$ непрерывна (следовательно, интегрируема) на $[a;\, b]$.

Обозначим $\rho_n = \sup\limits_{[a;\, b]} |f_n(x) - f(x)|$; по условию $\lim\limits_{n \to \infty} \rho_n = 0$.

Тогда при всех $x \in [a;\, b]$:
$$
\left|\int_a^x f_n(t)\,dt - \int_a^x f(t)\,dt\right| \le \int_a^x |f_n(t) - f(t)|\,dt \le \rho_n \cdot (x - a) \le \rho_n(b - a),
$$
откуда
$$
\sup_{[a;\, b]} \left|\int_a^x f_n(t)\,dt - \int_a^x f(t)\,dt\right| \le \rho_n(b - a) \to 0.
$$
\end{proof}

\begin{notion}
Равномерная сходимость является достаточным, но не необходимым условием возможности совершения предельного перехода под знаком интеграла.
\end{notion}

\section{Теорема о дифференцировании}

\begin{thm}[Почленное дифференцирование] \label{ch16:thm:diff}
Пусть функции $f_n$ непрерывно дифференцируемы на отрезке $[a;\, b]$, причём $f'_n(x) \rightrightarrows \varphi(x)$, $x \in [a;\, b]$, а последовательность $f_n$ сходится в некоторой точке $x_0 \in [a;\, b]$.

Тогда $f_n(x) \rightrightarrows f(x)$, $x \in [a;\, b]$, где функция $f$ непрерывно дифференцируема на $[a;\, b]$, причём
$$
f'(x) = \varphi(x) \quad \forall\, x \in [a;\, b],
$$
т.е. $\displaystyle\lim_{n \to \infty} f'_n(x) = \left(\lim_{n \to \infty} f_n(x)\right)'$ "--- предельный переход под знаком производной.

В концах отрезка все производные предполагаются односторонними.
\end{thm}

\begin{proof}
По теореме~\ref{ch16:thm:cont} функция $\varphi$ непрерывна на $[a;\, b]$, а по теореме~\ref{ch16:thm:int}:
$$
\int_a^x f'_n(t)\,dt \rightrightarrows \int_a^x \varphi(t)\,dt, \quad \text{т.е.} \quad f_n(x) - f_n(a) \rightrightarrows \int_a^x \varphi(t)\,dt, \quad x \in [a;\, b].
$$

При $x = x_0$:
$$
f_n(x_0) - f_n(a) \to \int_a^{x_0} \varphi(t)\,dt.
$$

Это соотношение не зависит от $x$, поэтому сходимость можно считать равномерной на $[a;\, b]$. Вычитая:
$$
\bigl(f_n(x) - f_n(a)\bigr) - \bigl(f_n(x_0) - f_n(a)\bigr) \rightrightarrows \int_a^x \varphi(t)\,dt - \int_a^{x_0} \varphi(t)\,dt,
$$
т.е.
$$
f_n(x) - f_n(x_0) \rightrightarrows \int_{x_0}^x \varphi(t)\,dt.
$$

По условию $\lim\limits_{n \to \infty} f_n(x_0) = A \in \bbR$. Тогда
$$
f_n(x) \rightrightarrows A + \int_{x_0}^x \varphi(t)\,dt =: f(x).
$$

Функция $f(x)$ непрерывно дифференцируема на $[a;\, b]$ и $f'(x) = \varphi(x)$ при всех $x \in [a;\, b]$.
\end{proof}

\begin{notion}
Если в теореме~\ref{ch16:thm:diff} отказаться от требования сходимости $f_n$ хотя бы в одной точке $x_0 \in [a;\, b]$, то даже при равномерной сходимости $f'_n$ последовательность $f_n$ может не сходиться ни в одной точке. Например, $f_n(x) = n$ не сходится ни в одной точке, а $f'_n(x) \equiv 0 \rightrightarrows 0$.
\end{notion}

\begin{notion}
Из равномерной сходимости последовательности непрерывно дифференцируемых функций к непрерывно дифференцируемой функции ещё не следует возможность предельного перехода под знаком производной.

Например, $f_n(x) = \dfrac{1}{n}\arctg x^n$, $x \in [-1;\, 1]$. Так как $|f_n(x)| \le \dfrac{\pi}{2n}$, имеем $f_n(x) \rightrightarrows 0 \equiv f(x)$, $f'(x) \equiv 0$. Но $f'_n(x) = \dfrac{x^{n-1}}{1 + x^{2n}}$, и $f'_n(1) = \dfrac{1}{2}$, $f'_n(-1) = \dfrac{(-1)^n}{2}$, так что $\lim\limits_{n \to \infty} f'_n(-1)$ не существует.
\end{notion}

\begin{notion}[Предостережение]
При исследовании равномерной сходимости функциональной последовательности нельзя заменять её общий член эквивалентной величиной. Отброшенные $o$-малые могут быть неравномерными по~$x$.

Например, $f_n(x) = \dfrac{1}{n} + \dfrac{x}{n^2}$, $\lim f_n(x) = 0$ при всех $x \in \bbR$. При всех $x$ имеем $f_n(x) \sim \dfrac{1}{n}$ при $n \to \infty$, и $\dfrac{1}{n} \rightrightarrows 0$. Однако $\sup\limits_{\bbR} |f_n(x)| = +\infty$ при всех $n$, поэтому $f_n(x)$ сходится к нулю неравномерно по $x \in \bbR$.
\end{notion}

\section{Равномерная сходимость функциональных рядов}

Члены числового ряда $\sum\limits_{n=1}^{\infty} u_n$ могут зависеть от параметра $x \in E$. Это значит, что на множестве $E$ определён \textit{функциональный ряд} $\sum\limits_{n=1}^{\infty} u_n(x)$.

\begin{defn}
Функциональный ряд $\sum\limits_{n=1}^{\infty} u_n(x)$ называется \textit{равномерно сходящимся} на множестве $E$, если последовательность его частичных сумм $S_n(x) = \sum\limits_{k=1}^{n} u_k(x)$ равномерно сходится на~$E$.
\end{defn}

Теоремы~\ref{ch16:thm:cont}--\ref{ch16:thm:diff} естественным образом переформулируются для функциональных рядов.

\begin{thmn}[Критерий Коши равномерной сходимости ряда] \label{ch16:thm:cauchy_series}
Функциональный ряд $\sum\limits_{n=1}^{\infty} u_n(x)$ равномерно сходится на множестве $E$ тогда и только тогда, когда
$$
\forall\, \varepsilon > 0 \;\; \exists\, n_0 \cquad \forall\, n \ge n_0, \;\; \forall\, p \in \bbN, \;\; \forall\, x \in E \cquad \left|\sum_{k=n+1}^{n+p} u_k(x)\right| < \varepsilon.
$$
\end{thmn}

\begin{thmn}[Непрерывность суммы ряда] \label{ch16:thm:cont_series}
Если функции $u_n$, $n \in \bbN$, непрерывны на множестве $E$ и ряд $\sum\limits_{n=1}^{\infty} u_n(x)$ равномерно сходится на $E$, то его сумма $S(x)$ непрерывна на~$E$.
\end{thmn}

\begin{thmn}[Почленное интегрирование ряда] \label{ch16:thm:int_series}
Если функции $u_n$, $n \in \bbN$, непрерывны на отрезке $[a;\, b]$ и ряд $\sum\limits_{n=1}^{\infty} u_n(x)$ равномерно сходится к сумме $S(x)$ на $[a;\, b]$, то ряд $\sum\limits_{n=1}^{\infty} \int_a^x u_n(t)\,dt$ равномерно сходится на $[a;\, b]$ к сумме $\int_a^x S(t)\,dt$.

В частности, при $x = b$:
$$
\sum_{n=1}^{\infty} \int_a^b u_n(t)\,dt = \int_a^b \left(\sum_{n=1}^{\infty} u_n(t)\right) dt.
$$
\end{thmn}

\begin{thmn}[Почленное дифференцирование ряда] \label{ch16:thm:diff_series}
Если функции $u_n$, $n \in \bbN$, непрерывно дифференцируемы на отрезке $[a;\, b]$, ряд $\sum\limits_{n=1}^{\infty} u'_n(x)$ равномерно на $[a;\, b]$ сходится к сумме $\varphi(x)$, а ряд $\sum\limits_{n=1}^{\infty} u_n(x)$ сходится в некоторой точке $x_0 \in [a;\, b]$, то $\sum\limits_{n=1}^{\infty} u_n(x)$ равномерно на $[a;\, b]$ сходится к непрерывно дифференцируемой функции $S(x)$, причём
$$
S'(x) = \varphi(x) \quad \forall\, x \in [a;\, b], \quad \text{т.е.} \quad \left(\sum_{n=1}^{\infty} u_n(x)\right)' = \sum_{n=1}^{\infty} u'_n(x).
$$
\end{thmn}